\subsection{A product of centered complex exponentials}
\label{sec:product-of-centered-exponentials}

This subsection requires only $\E X=0$ and $\E|X|<\infty$.
Neither a density nor a second moment is needed.
Write
$\phi(u)=\E e^{iuX}$, and for independent copies $X_1,X_2$ put
\[
 O_\theta=\begin{pmatrix}\cos\theta&-\sin\theta\\
                         \sin\theta&\cos\theta\end{pmatrix},
 \qquad
 W_{\xi,\eta}(x,y)
   =(e^{i\xi x}-\phi(\xi))(e^{i\eta y}-\phi(\eta)).
\]
This fixes the sign convention for rotation.
Centering removes the two marginal responses.
It leaves only the dependence created by rotation.
In particular,
\begin{equation}
 \E W_{\xi,\eta}(X_1,X_2)=0,
 \qquad
 \E|W_{\xi,\eta}(X_1,X_2)|^2
 =(1-|\phi(\xi)|^2)(1-|\phi(\eta)|^2)\le1.
 \label{eq:centered-exponential-product-variance}
\end{equation}
Pointwise, $|W_{\xi,\eta}|\le4$, so its variance under an arbitrary
comparison law is at most $16$.  Both bounds are dimension-free and require
no tail assumption.

\subsection{Gaussian characterization by a global linear ODE}
\label{sec:fourier-ode}

The next lemma makes the rotational Fourier defect a logarithm-free rigidity
tool, valid even when \(\phi\) has zeros.  Contrapositively, every
non-Gaussian law has a nonzero rotational response.

\begin{lemma}[Gaussian characterization by the rotational Fourier defect]
\label{lem:fourier-ode}
For the product of centered complex exponentials defined above,
\begin{equation}
 R(\xi,\eta)
 =\left.\frac{d}{d\theta}
       \E W_{\xi,\eta}(O_\theta(X_1,X_2))\right|_{\theta=0}
   =\eta\phi'(\xi)\phi(\eta)
          -\xi\phi(\xi)\phi'(\eta).
 \label{eq:rotational-fourier-response}
\end{equation}
Moreover,
\begin{equation}
 R\equiv0
 \quad\Longleftrightarrow\quad
 \Law(X)\text{ is a centered Gaussian, possibly degenerate}.
 \label{eq:fourier-gaussian-characterization}
\end{equation}
Equivalently, all pairs of values of the weighted Fourier jet
\[
 u\longmapsto\bigl(u\phi(u),\phi'(u)\bigr)
\]
are linearly dependent if and only if \(\Law(X)\) is a centered Gaussian,
possibly degenerate.
\end{lemma}
\begin{proof}
The first moment implies $\phi\in C^1(\R)$.
The joint characteristic function after rotation is
\[
 \phi(\xi\cos\theta+\eta\sin\theta)
 \phi(-\xi\sin\theta+\eta\cos\theta).
\]
The marginal characteristic functions are
$\phi(\xi\cos\theta)\phi(-\xi\sin\theta)$ and
$\phi(\eta\sin\theta)\phi(\eta\cos\theta)$.
Their derivatives at zero vanish because $\phi'(0)=i\E X=0$.
Differentiating the joint characteristic function therefore gives
\eqref{eq:rotational-fourier-response}.

If $R$ vanishes everywhere, continuity and $\phi(0)=1$ allow us to
choose $\eta_0\ne0$ with $\phi(\eta_0)\ne0$.
Fix this single frequency.
The defect equation then becomes the global ODE
\[
 \phi'(\xi)=c\xi\phi(\xi),\qquad
 c=\frac{\phi'(\eta_0)}{\eta_0\phi(\eta_0)},\qquad
 \phi(0)=1.
\]
No division by $\phi(\xi)$ has occurred.
Uniqueness for this linear ODE gives $\phi(\xi)=\exp(c\xi^2/2)$ on all of
$\R$, including any putative zero of $\phi$.
The identity $\phi(-\xi)=\overline{\phi(\xi)}$ forces $c\in\R$.
Also $|\phi(\xi)|\le1$ forces $c\le0$.
Write $c=-\sigma^2$.
The Gaussian conclusion follows.
If $\sigma=0$, we get the point mass at zero.
Substitution proves the converse.
\end{proof}

\subsection{The determinant representation and finite detection}
\label{sec:fourier-determinant}

The determinant form supplies the finite frequency pair used to construct the
dual moments in \cref{sec:fourier-dual-moments}.

Consider the curve
\[
 u\longmapsto\bigl(u\phi(u),\phi'(u)\bigr)\in\mathbb C^2,
\]
described as the weighted Fourier jet.  Its values satisfy
\begin{equation}
 \det\!\begin{pmatrix}
 \xi\phi(\xi)&\phi'(\xi)\\
 \eta\phi(\eta)&\phi'(\eta)
 \end{pmatrix}
 =-R(\xi,\eta).
 \label{eq:fourier-wedge}
\end{equation}
The weighting by \(u\) is essential: for a nondegenerate Gaussian law, the
unweighted curve \(u\mapsto(\phi(u),\phi'(u))\) is not contained in a fixed
line.

For a non-Gaussian law, continuity supplies distinct frequencies and a compact
detecting neighborhood, independent of the ambient dimension; the diagonal
\(\xi=\eta\) never detects the law.

The same product of centered complex exponentials gives a local Hellinger
lower bound for one rotation.
Choose a phase $\alpha$ with
$\operatorname{Re}(e^{-i\alpha}R(\xi,\eta))>0$ and apply
\[
 |\E_\lambda F-\E_\nu F|
 \le h(\lambda,\nu)
             \{2\E_\lambda F^2+2\E_\nu F^2\}^{1/2}
\]
to $F=\operatorname{Re}(e^{-i\alpha}W_{\xi,\eta})$.
It follows that
$h((O_\theta)_\#\Law(X)^{\otimes2},
\Law(X)^{\otimes2})\ge c|\theta|$ for small
$|\theta|$.
A first moment suffices for this lower bound.
The dimension-free matrix estimates below require the additional
second-moment and Sobolev hypotheses.

\subsection{Dimension-free variance and affine normalization}
\label{sec:variance-and-affine-normalization}

The dimension-free variance bounds needed below are already contained in
\eqref{eq:centered-exponential-product-variance}; we record only how \(R\)
changes under affine normalization.

\paragraph{Centering, scaling, and phase.}
Let $Y$ have mean \(m\) and characteristic function \(\phi_Y\), and put
\(X=(Y-m)/\sigma\), where \(\sigma>0\).
For \(a=\xi/\sigma\) and \(b=\eta/\sigma\), the rotational Fourier defect
of $X$ satisfies
\begin{equation}
 \begin{aligned}
 R(\xi,\eta)
 &=e^{-i(m/\sigma)(\xi+\eta)}
 \bigl[b\phi_Y'(a)\phi_Y(b)-a\phi_Y(a)\phi_Y'(b)\\
 &\hspace{8em}+im(a-b)\phi_Y(a)\phi_Y(b)\bigr].
 \end{aligned}
 \label{eq:centering-phase}
\end{equation}
The bracketed term is the mean-corrected rotational response for two copies of
\(Y\).  Centering contributes the displayed phase, scaling changes the
frequencies, and the modulus at corresponding frequencies is unchanged.  For
a centered symmetric law the defect is real, but its phase at a single pair is
not a complete measure of skewness.

For the concrete comparison $Y\sim\operatorname{Gamma}(2,1)$,
$m=2$, $\sigma=\sqrt2$, and $(\xi,\eta)=(1,2)$, the expression in brackets in
\eqref{eq:centering-phase} equals $4/27$, so the law centered with variance one has
\[
 R(1,2)=\frac4{27}e^{-3i\sqrt2}.
\]
Thus the defect can remain explicit even for an asymmetric law with a support
boundary.
